\documentclass[english,twoside, openright, hidelinks, 11 pt, class=report,crop=false]{standalone}
\input{../../preamb}
\input{../bm}

\begin{document}
\section{\brgrpr \label{brgrpr}}
\reg[\brdef \label{brdef}]{
	En brøk\index{brøk} er en annen måte å skrive et delestykke på. I en brøk kaller vi dividenden for \outl{teller}\index{teller} og divisoren for \outl{nevner}\index{nevner}.
	\begin{figure}
		\centering
		\includegraphics[]{\figp{br0_bm}}
	\end{figure}
}\vsk

\spr{
Vanlige måter å si $ \frac{1}{4} $ på er\footnote{I tillegg har vi utsagnene fra språkboksen på side \pageref{sprakdiv}.}
\begin{itemize}
	\item ''én firedel''
	\item ''1 av 4''
	\item ''1 over 4''
\end{itemize}
}
\newpage
\subsubsection{Brøk som mengde}
La oss se på brøken $ \frac{1}{4} $ som en mengde. Vi starter med å definere en rute\footnote{Her har vi valgt en enerrute som er større enn den vi brukte i \refkap{Talavare}.} med verdi 1: 
\begin{figure}[hbt]
	\centering
	\includegraphics{\figp{br1}}
\end{figure}
Så  deler vi denne ruten inn i fire mindre ruter som er like store. Hver av disse rutene er da $ \frac{1}{4} $ (1 av 4):
\begin{figure}[hbt]
	\centering
	\includegraphics{\figp{br2}}
\end{figure} 

Har vi én slik rute, har vi altså 1 firedel:
\begin{figure}[hbt]
	\centering
	\includegraphics{\figp{br3a}}
\end{figure} 
Men skal man fra en figur kunne se hvor stor en brøk er, må man vite hvor stor 1 er, og for å få dette lettere til syne skal vi også ta med de ''tomme'' rutene:
\fig{br3}
\newpage
Slik vil de blå og de tomme rutene fortelle oss hvor mange biter 1 er delt inn i, mens de blå rutene alene forteller oss hvor mange sånne biter det \textsl{egentlig} er. Altså kan vi si at\regv
 
\st{ \vs
\alg{
	\text{antall blå ruter}&=\text{teller} \\	
	\text{antall blå ruter}+\text{antall tommme ruter}&=\text{nevner} 
}
\fig{br0a}
}

\subsubsection{Brøk på tallinja}
På tallinja deler vi lengden mellom 0 og 1 inn i like mange lengder som nevneren angir. Har vi en brøk med 4 i nevner, deler vi lengden mellom 0 og 1 inn i 4 like lengder:
\fig{br9}
Tallinja er fin å bruke for å tegne brøker som er større enn 1:
\fig{br9a}
\subsubsection{Teller og nevner oppsummert}
Selv om vi har vært innom det allerede, er det så avgjørende å forstå hva telleren og nevneren sier oss at vi tar en kort oppsummering:\regv
\st{
\begin{itemize}
	\item Nevneren forteller hvor mange like store biter 1 \\ er delt inn i.
	\item Telleren forteller hvor mange slike biter det er.
\end{itemize}
}
\newpage
\section{\brvu}
\reg[Verdien til en brøk]{Verdien til en brøk finner vi ved å dele telleren med nevneren.}
\eks{
Finn verdien til $ \dfrac{1}{4} $.

\sv 
\[ \frac{1}{4}=0,25 \]
(Se \refkap{Utrekning} for hvordan man kan finne at $ 1:4=0,25 $.)
}
\subsubsection{Brøker med lik verdi}
Brøker kan ha lik verdi selv om de ser forskjellige ut. De kalles da \outl{likeverdige}\index{likeverdig} brøker. Hvis du regner ut $ 1:2 $, $ 2:4 $ og $ 4:8 $, får du i alle tilfeller 0,5 som svar. Dette betyr at
\alg{
\frac{1}{2}=\frac{2}{4}=\frac{4}{8}=0,5
} \\[5pt]
\amounts{\fig{br12}}
\numrline{\fig{br12a}}
\subsubsection{Utviding}
At brøker kan se forskjellige ut, men ha lik verdi, betyr at vi kan\\ endre på utseendet til en brøk uten å endre verdien. La oss som \\eksempel gjøre om $ \frac{3}{5} $ til en brøk med lik verdi, men med 10 som nevner:
\begin{itemize}
	\item $ \frac{3}{5} $ kan vi gjøre om til en brøk med 10 i nevner om vi deler hver femdel inn i 2 like biter, for da er 1 til sammen delt inn i\\ $ {5\cdot2=10} $ biter.
	\item Telleren i $ \frac{3}{5} $ viser at der er 3 femdeler. Når disse deles i to, blir de totalt til $ 3\cdot2=6 $ tideler. Altså har $ \frac{3}{5} $ lik verdi som $ \frac{6}{10} $.
\end{itemize}
\[ \frac{3}{5}=\frac{3\cdot2}{5\cdot2}=\frac{6}{10} \]
\amounts{\fig{br13}}
\numrline{\fig{br13a}}
\subsubsection{Forkorting}
Siden vi kan skrive om $\frac{3}{5}$ til $\frac{6}{10}$ ved å gange både teller og nevner med 2, må vi også kunne gjøre det omvendte. $ \frac{6}{10} $ kan vi skrive om til $\frac{3}{5}$ ved å dele både teller og nevner med 2:
\[ \frac{6}{10}=\frac{6:2}{10:2}=\frac{3}{5} \]
\newpage
\reg[Utviding og forkorting av brøk]{
Vi kan gange eller dele teller og nevner med det samme tallet uten at brøken endrer verdi. \vsk

Å gange med et tall større enn 1 kalles å \outl{utvide}\index{brøk!utviding av} brøken. Å dele med et tall større enn 1 kalles å \outl{forkorte}\index{brøk!forkorting av} brøken.
}
\eks[1]{
Utvid $\frac{7}{3}$ med 4.

\sv
\[ \frac{7}{3}=\frac{7\cdot4}{3\cdot4}=\frac{21}{12} \]
}
\eks[2]{Utvid $ \frac{3}{5} $ til en brøk med 20 som nevner.
	
	\sv
	Da $ {5\cdot4=20} $, utvider vi brøken med 4:
	\[ \frac{3}{5} = \frac{3\cdot4}{5\cdot4} = \frac{12}{20} \]
}
\eks[3]{
Forkort $\frac{6}{10}$ med 2.

\sv
\[ \frac{6}{10}=\frac{6:2}{10:2}=\frac{3}{5} \]
}
\newpage
\eks[4]{
	Forkort $ \frac{18}{30} $ til en brøk med 5 som nevner.
	
	\sv
	Da $ 30:6=5 $, forkorter vi brøken med 6:
	\[ \frac{18}{30}=\frac{18:6}{30:6} =\frac{3}{5} \]
}
\eks[5]{
Undersøk om brøkene $\frac{2}{3}$ og $\frac{5}{7}$ har lik eller ulik verdi.

\sv
Vi utvider begge brøkene slik at de har lik nevner. Ved å utvide $\frac{2}{3}$ med 7 og $\frac{5}{7}$ med 3, får vi at
\[ \frac{2}{3}=\frac{2\cdot7}{3\cdot7}=\frac{14}{21}\qquad,\qquad\frac{5}{7}=\frac{5\cdot3}{7\cdot3}=\frac{15}{21} \]
Vi har at
\[ \frac{14}{21}<\frac{15}{21} \]
Altså er
\[ \frac{2}{3}<\frac{5}{7} \]
}
\newpage
\section{\bradsub}
Addisjon og subtraksjon av brøker handler i stor grad om nevnerne. Husk at nevnerne viser inndelingen av 1. Hvis brøker har lik nevner, representerer de et antall biter med lik størrelse. Da gir det mening å regne addisjon eller subtraksjon mellom tellerne. Hvis brøker har ulike nevnere, representerer de et antall biter med ulik størrelse, og da gir ikke addisjon eller subtraksjon mellom tellerne direkte mening.
\subsubsection{Lik nevner}
Om vi for eksempel har 2 seksdeler og adderer 3 seksdeler, ender vi opp med 5 seksdeler:
\[ \frac{2}{6}+\frac{3}{6}=\frac{5}{6} \]
\amounts{\fig{br4}}
\numrline{\fig{br4a}}
\reg[Addisjon/subtraksjon av brøker med lik nevner \label{bradliknemn}]{
Når vi regner addisjon/subtraksjon mellom brøker med lik nevner, finner vi summen/differansen av tellerne, og beholder nevneren.
}
\eks[1]{ \vs
\[ \frac{2}{7}+\frac{8}{7}=\frac{2+8}{7} = \frac{10}{7} \]
}
\eks[2]{ \vs
\[ \frac{7}{9}-\frac{5}{9}=\frac{7-5}{9} = \frac{2}{9} \]
}
\subsection*{Ulike nevnere}
La oss se på regnestykket\footnote{Vi minner om at rødfargen på pila viser at man skal vandre fra pilspissen til den andre enden.} \vspace{-5pt}
\[ \frac{3}{5}-\frac{1}{2} \]
\amounts{\fig{br7} \vs} \vs
\numrline{\fig{br7t} \vs}
Skal vi skrive differansen som en brøk, må vi sørge for at brøkene har lik nevner. De to brøkene våre kan begge ha 10 som nevner:
\alg{
\frac{3}{5}=\frac{3\cdot2}{5\cdot2}=\frac{6}{10}\quad\quad,\quad\quad \frac{1}{2}=\frac{1\cdot5}{2\cdot5}=\frac{5}{10}
}
Dette betyr at
\[ \frac{3}{5}-\frac{1}{2}=\frac{6}{10}-\frac{5}{10} \]
\amounts{\fig{br7a}} \vs
\numrline{\fig{br7ta} }
Det vi har gjort, er å utvide begge brøkene slik at de har lik nevner, nemlig 10. Når nevnerne i brøkene er like, kan vi regne ut differansen mellom tellerne:
\[ \frac{3}{5}-\frac{1}{2}=\frac{6}{10}-\frac{5}{10}=\frac{1}{10} \]
\reg[Addisjon/subtraksjon av brøker med ulik\\ nevner \label{braduliknemn}]{
Når vi regner addisjon/subtraksjon mellom brøker med ulik nevner, må vi utvide brøkene slik at de har lik nevner, for så å bruke \rref{bradliknemn}.
}
\eks[1]{ \label{eksadbrok}
Regn ut
\[ \frac{2}{\colb{9}}+\frac{6}{\colo{7}} \]
\sv \vs
Begge nevnerne kan bli $ 63 $ hvis vi ganger med rett heltall. Vi utvider derfor til brøker med 63 i nevner:
\[ \frac{2\cdot\colo{7}}{9\cdot \colo{7}}+\frac{6\cdot\colb{9}}{7\cdot\colb{9}}=\frac{14}{63}+\frac{54}{63}=\frac{68}{63} \]
}

\eks[2]{
Regn ut 
\[ \frac{2}{5}+9 \]

\sv 
9 kan vi skrive som $\frac{9}{1}$. Altså er
\[ \frac{2}{5}+9=\frac{2}{5}+\frac{9}{1} \]
Hvis vi utvider $\frac{9}{1}$ med 5, får vi to brøker med lik nevner:
\[ \frac{2}{5}+\frac{9\cdot5}{1\cdot5}=\frac{2}{5}+\frac{45}{5}=\frac{47}{5} \]
}
\newpage
\info{Fellesnevner}{
I \textsl{Eksempel 1} på side \pageref{eksadbrok} blir 63 kalt en \outl{fellesnevner}\index{fellesnevner}. Dette fordi det finnes heltall vi kan gange nevnerne med som gir oss tallet 63:
\algv{
	9\cdot7 &= 63 \\
	7\cdot 9 &= 63 
}
Hvis vi ganger sammen alle nevnerne i et regnestykke, finner vi alltid en fellesnevner, men vi sparer oss for store tall om vi finner den \textsl{minste} fellesnevneren. Ta for eksempel regnestykket
\[ \frac{7}{6}+\frac{5}{3} \]
Her kan vi bruke fellesnevneren $ {6\cdot3=18} $, men det er bedre å merke seg at $ 6\cdot1=3\cdot2=6 $ også er en fellesnevner. Altså er
\alg{
	\frac{7}{6}+\frac{5}{3}&=\frac{7}{6}+\frac{5\cdot2}{3\cdot2}\br
	&=\frac{7}{6}+\frac{10}{6}\br
	&=\frac{17}{6}
}

}
\eks[3]{
	Regn ut
	\[ \frac{3}{2}-\frac{5}{8}+\frac{10}{4} \]
	
	\sv
	Alle nevnerne kan bli 8 hvis vi ganger med rett heltall. Vi utvider derfor brøkene slik at de har 8 som nevner:
	\alg{
		\frac{3}{2}-\frac{5}{8}+\frac{10}{4} &= \frac{3\cdot4}{2\cdot4}-\frac{5}{8}+\frac{10\cdot2}{4\cdot2} \br
		&= \frac{12}{8}-\frac{5}{8}+\frac{20}{8} \br
		&= \frac{27}{8}
	}
}
\newpage

\section{\brgngheil}
I \hrs{Gonging}{seksjon} så vi at ganging med heltall er det samme som gjentatt addisjon. Skal vi for eksempel regne ut $ \frac{2}{5}\cdot 3 $, kan vi regne slik:
\alg{
\frac{2}{5}\cdot 3 = \frac{2}{5}+\frac{2}{5}+\frac{2}{5}
= \frac{2+2+2}{5} = \frac{6}{5}
}
\numrline{\fig{br14}}
Men vi vet også at \y{2+2+2=2\cdot3}, og derfor kan vi forenkle \\ regnestykket vårt:
\[ \frac{2}{5}\cdot 3 = \frac{2\cdot3}{5}= \frac{6}{5} \]
Multiplikasjon mellom heltall og brøk er også kommutativ\footnote{Husk at $ \frac{2}{5} $ bare er en omskriving av $ {2:5} $, og at vi av \rref{rret} regner fra venstre mot høyre. }:
\alg{
3\cdot\frac{2}{5}&=3\cdot2:5 \\
&=6:5 \\
&=\frac{6}{5}
}
\reg[\label{brhel}Brøk ganget med heltal]{
Når vi ganger en brøk med et heltall, ganger vi heltallet med telleren i brøken.
}
\eks[1]{
\[ \frac{1}{3}\cdot 4=\frac{1\cdot 4}{3} \br =\frac{4}{3} \]
}
\eks[2]{
\[ 3\cdot\frac{2}{5}=\frac{3\cdot 2}{5} \br =\frac{6}{5} \]
}


\newpage
\section{\brdelheil} \label{brdelheil}
Tidligere i boka har vi sett at
\begin{itemize}
	\item Deling kan man se på som en lik fordeling av en mengde.
	\item I en brøk er det telleren som forteller noe om antallet vi har av en mengde (nevneren forteller om inndelingen av 1).
\end{itemize}
Når en brøk skal deles med et heltall, er det dermed telleren som gir antallet som skal fordeles likt.
\subsubsection{Tilfellet der telleren er delelig med divisoren}
La oss regne ut
\[ \frac{6}{8}:2 \]
Vi har da 6 åttedeler som vi skal fordele likt på 2 grupper. $ 6:2=3 $, altså får vi 3 åttedeler i hver gruppe. Dermed er
\[ \frac{6}{8}:2=\frac{3}{8} \]
\amounts{\fig{br15a}}
\numrline{\fig{br15ta} \vs}
\newpage
\subsubsection{Tilfellet der telleren ikke er delelig med divisoren}
La oss regne ut
\[ \frac{3}{4}:2 \]
Vi kan alltid utvide brøken vår slik at telleren blir delelig med divisoren. Siden vi skal dele med 2, utvider vi brøken vår med 2:
\[ \frac{3}{4}=\frac{3\cdot2}{4\cdot2}=\frac{6}{8} \]
Vi har akkurat sett at ${\frac{6}{8}:2=\frac{3}{8}}$. Da ${\frac{3}{4}=\frac{6}{8}}$, har vi at
\[ \frac{3}{4}:2=\frac{6}{8}:2=\frac{3}{8} \]
Rent matematisk har vi ganget nevneren til $ \frac{3}{4} $ med 2:
\[ \frac{3}{4}:2 = \frac{3}{4\cdot2} = \frac{3}{8} \]
\reg[Brøk delt med heltall \label{brdeenk}]{
Når vi deler en brøk med et heltall, ganger vi nevneren med heltallet.
}
\eks[1]{\vs
\[ \frac{5}{3}:6 = \frac{5}{3\cdot 6} 
= \frac{5}{18} \]
}
\info{\note}{
På side \pageref{brdelheil} fant vi at
\[ \frac{6}{8}:2=\frac{3}{8} \]
Da ganget vi ikke nevneren med 2 slik \rref{brdeenk} tilsier. Om vi gjør det, får vi 
\alg{
\frac{6}{8}:2=\frac{6}{8\cdot2}=\frac{6}{16}
}
$\frac{6}{16}$ forkortet med 2 er lik $\frac{3}{8}$, de to svarene har altså (selvsagt) lik verdi. Skal vi dele en brøk på et heltall, og telleren er delelig med heltallet, kan vi direkte dele telleren på heltallet. I slike tilfeller er det altså ikke \textsl{feil}, men heller \textsl{ikke nødvendig} å bruke \rref{brdeenk}.
}
\section{\brgngbr \label{brgngbr}}		
La oss regne ut
\[  {\frac{5}{4}\cdot\frac{3}{2}}\] 
Dette kan vi skrive som
\[ \frac{5}{4}\cdot3:2 \]
Da skal vi først gange $ \frac{5}{4} $ med 3, og så dele produktet med 2. Av \rref{brhel} er
\alg{
\frac{5}{4}\cdot3 =\frac{5\cdot 3}{4}
}
Og av \rref{brdeenk} er
\alg{
\frac{5\cdot3}{4}:2=\frac{5\cdot3}{4\cdot2}
}
Altså er 
\alg{
\frac{5}{4}\cdot \frac{3}{2}=\frac{5\cdot 3}{4\cdot2}
}
\reg[\brtbr\label{brtbr}]{
Når vi ganger to brøker med hverandre, ganger vi teller med teller og nevner med nevner.
}
\eks[1]{ \vs
\[ \frac{4}{7}\cdot \frac{6}{9}= \frac{4\cdot6}{7\cdot9}= \frac{24}{63} \]
}
\eks[2]{ \vs
\[ \frac{1}{2}\cdot\frac{9}{10}=\frac{1\cdot9}{2\cdot10}=\frac{9}{20} \]
}
\newpage
\section{\brkans}
Når telleren og nevneren har lik verdi, er verdien til brøken alltid 1. For eksempel er \y{\frac{3}{3}=1}, \y{\frac{25}{25}=1} og så videre. Dette kan vi utnytte for å forenkle brøkuttrykk. \vsk

La oss forenkle brøkuttrykket
\[ \frac{8\cdot 5}{9\cdot8} \]
Da $ {8\cdot5}={5\cdot8} $, kan vi skrive
\[ \frac{8\cdot5}{9\cdot8}=\frac{5\cdot8}{9\cdot8} \]
Og som vi nylig har sett (\rref{brtbr}) er
\[\frac{5\cdot8}{9\cdot8}= \frac{5}{9}\cdot\frac{8}{8}\]
Siden $ {\frac{8}{8}=1} $, har vi at
\[ \frac{5}{9}\cdot\frac{8}{8} = \frac{5}{9}\cdot1 = \frac{5}{9} \]
Når bare består av faktorer, kan man alltid omrokkere slik vi har gjort over, men når man har forstått hva omrokkeringen ender med, er det bedre å bruke \outl{kansellering}\footnote{I \refkap{Rekneartane} brukte vi ordet \textit{kansellering} om noe annet. Ordet brukes forskjellig ved kombinasjonen av addisjon/subtraksjon og av multiplikasjon/divisjon.}\index{kansellering}. Man setter da en strek over to og to like faktorer for å vise at de utgjør en brøk med verdien 1. Tilfellet vi akkurat så på skriver vi da som 
\[ \frac{\cancel{8}\cdot 5}{9\cdot\cancel{8}}=\frac{5}{9} \]
\newpage
\reg[Kansellering av faktorer\label{kans}]{Når en brøk bare består av faktorer, kan vi kansellere par av like faktorer i teller og nevner.}
\eks[1]{
	Kanseller så mange faktorer som mulig i brøken
	\[ \frac{3\cdot12\cdot7}{7\cdot 4 \cdot 12} \]
	\sv \vsb
\[\frac{3\cdot\cancel{12}\cdot\cancel{7}}{\cancel{7}\cdot 4 \cdot \cancel{12}} = \frac{3}{4}   \]
}
\eks[2]{ \label{eksbrkan}
	Forkort brøken $ \frac{12}{42} $.
	
	\sv 
	$ 6 $ er en faktor i både $ 12 $ og $ 42 $, altså er	
	\alg{
		\frac{12}{42} = \frac{\cancel{6}\cdot2}{\cancel{6}\cdot7}= \frac{2}{7}
	}
}
\eks[3]{
	Forkort brøken $ \frac{48}{16}  $.
	
	\sv
	$ 16 $ er en faktor i $ 48 $, altså er
	\alg{
		\frac{48}{16}&=
		\frac{3\cdot\cancel{16}}{\cancel{16}}\br
		&= \frac{3}{1}\br
		&= 3
	}
	\mers{Hvis alle faktorer er kansellert i telleren/nevneren, er telleren/nevneren lik 1.}
}
\newpage
\info{Forkorting ved primtalsfaktorisering}{
	Det er ikke alltid like lett å legge merke til en felles faktor, slik vi har gjort i \textsl{Eksempel 2} og \textsl{Eksempel 3} på side \pageref{eksbrkan}. Vil man være helt sikker på at man ikke har oversett felles faktorer, kan man primtalsfaktorisere (se \refsec{Faktorisering}) både teller og nevner. For eksempel har vi at
\[ \frac{12}{42}=\frac{\cancel{2}\cdot2\cdot\cancel{3}}{\cancel{2}\cdot\cancel{3}\cdot7}=\frac{2}{7} \]
}
\info{Brøker forenkler utregninger}{
	Desimaltallet $ 0,125 $ kan vi skrive som brøken $ \frac{1}{8} $. Et regnestykke som ${0,125\cdot16}$ kan vi da forenkle slik:
	\alg{
		0,125\cdot16 &= \frac{1}{8}\cdot16 \br
		&= \frac{2 \cdot\cancel{8}}{\cancel{8}}\\
		&= 2
	}
} 
\info{''Å stryke nuller''}{
	Et tall som 3000 kan vi skrive som $ 3\cdot10\cdot10\cdot10 $, mens 700 kan vi skrive som $ 7\cdot10\cdot10 $. Brøken $ \frac{3000}{700} $ kan vi derfor forkorte slik:
	\alg{
		\frac{3000}{700}&= \frac{3\cdot\cancel{10}\cdot\cancel{10}\cdot10}{7\cdot\cancel{10}\cdot\cancel{10}}\br
		&= \frac{3\cdot10}{7} \br
		&= \frac{30}{7}
	}
	I praksis er dette det samme som ''å stryke nuller'':
	\alg{
		\frac{30\cancel{00}}{7\cancel{00}}&= \frac{30}{7}
	}
	\textsl{Obs!} Null er de eneste sifferet vi kan ''stryke'' slik, for eksempel kan vi ikke forkorte $ \frac{123}{13} $ på noen som helst måte. I tillegg kan vi bare ''stryke'' nuller som står som bakerste siffer, for eksempel kan vi ikke ''stryke'' nuller i brøken $ \frac{101}{10} $.  
}
\newpage
\section{\brdelmbr}
\subsubsection{Deling ved å se på tallinja}
La oss regne ut $ {4:\frac{2}{3} } $. Siden brøken vi deler 4 på har 3 i nevner, kan det være en idé å gjøre om også 4 til en brøk med $ 3 $ i nevner. Vi har at 
\[ 4=\dfrac{12}{3} \]
\numrline{\fig{br10c4}}
I \refkap{Rekneartane} så vi at utregningen av $ 4:\frac{2}{3} $ er svaret på spørsmålet
\begin{center}
	''Hvor mange ganger $ \frac{2}{3} $ går på 4''
\end{center}
Ved å se på tallinja, finner vi at $ \frac{2}{3} $ går 6 ganger på 4. Altså er
\[ 4:\frac{2}{3}=6 \]
\numrline{\fig{br10c}}
\newpage
\subsubsection{En generell metode}
Vi kan ikke se på en tallinje hver gang vi skal dele med brøker, så nå skal vi komme fram til en generell regnemetode ved igjen å bruke $ 4:\frac{2}{3} $ som eksempel. For denne metoden bruker vi denne betydningen av divisjon:
\begin{center}
	$  4:\dfrac{2}{3}= $ ''Tallet vi må gange $ \frac{2}{3} $ med for å få 4''
\end{center}
For å finne dette tallet starter vi med å gange $ \frac{2}{3} $ med tallet som gjør at produktet er 1. Dette er \outl{den omvendte brøken}\index{brøk!omvendt} av $ \frac{2}{3} $, som er $ \frac{3}{2} $:
\[ \frac{2}{3}\cdot\frac{3}{2}=1 \]
Nå gjenstår det bare å gange med 4 for å få 4:
\[ \frac{2}{3}\cdot\frac{3}{2}\cdot4=4 \]
$ \frac{3}{2}\cdot4 $ er altså tallet vi må gange $ \frac{2}{3} $ med for å få 4 . Dette betyr at
\alg{
4:\frac{2}{3}&=\frac{3}{2}\cdot4 \br
&=\frac{12}{2}\br
&=6
}
\newpage
\reg[\delmbr \label{delmbr}]{
Når vi deler et tall med en brøk, ganger vi tallet med den omvendte brøken.
}
\eks[1]{
\[ 3:\frac{\color{c1}2}{\color{c2}9}=3\cdot\frac{\color{c2}9}{\color{c1}2}=\frac{27}{2} \]
}
\eks[2]{
\[ \frac{4}{3}:\frac{\color{c1}5}{\color{c2}8} = \frac{4}{3}\cdot\frac{\color{c2}8}{\color{c1}5}	= \frac{32}{15} \]
}
\eks[3]{
\[ 		\frac{3}{5}:\frac{3}{10}= \frac{3}{5}\cdot\frac{10}{3} = \frac{30}{15} \]
	Her bør vi også se at brøken kan forkortes:
\[ \frac{30}{15}= \frac{2\cdot\cancel{15}}{ \cancel{15}} = 2 \]
\mer Vi kan spare oss for store tall hvis vi kansellerer faktorer underveis i utregninger:
\[ \frac{3}{5}\cdot\frac{10}{3} = \frac{\cancel{3}\cdot2\cdot\cancel{5}}{\cancel{5}\cdot\cancel{3}} = 2 \]
}
\newpage
\section{\Rasjtal \label{Rasjtal}}
\reg[Rasjonale tall]{
Et tall som kan bli skrevet som en brøk med heltalls teller og nevner, er et \outl{rasjonalt tall}\index{tall!rasjonalt}.
}
\info{Merk}{Rasjonale tall gir oss en samlebetegnelse for
\begin{itemize}
	\item \textbf{Heltall} \\
	For eksempel $ 4=\frac{4}{1} $.
	\item \textbf{Desimaltall med endelig antall desimaler}\\
	For eksempel $ 0,2=\frac{1}{5} $.
	\item \textbf{Desimaltall med repeterende desimalmønster} \\
	For eksempel\footnote{\sym{$ \bar{3} $} viser at 3 fortsetter i det uendelige. En annen måte å vise dette på er å bruke symbolet \sym{...}\,. Altså er $ 0,08\bar{3}=0,08333333 ... $} $ 0,08\bar{3}=\frac{1}{12} $. 
\end{itemize}} \vsk

\reg[Blandet tall\index{tall!blandet}]{
	Om vi adderer et heltall med en brøk der telleren er mindre enn nevneren, får vi et \outl{blandet tall}.
}
\eks[1]{
	Tre forskjellige blandede tall:
	\[   2+\dfrac{5}{7} \qquad\qquad
	8+\dfrac{2}{7} \qquad\qquad
	\dfrac{1}{10}+4  \]
}
\info{Obs!}{
	I mange tekster vil du finne tall som dem fra \textsl{Eksempel 1} skrevet slik:
	\[   2\dfrac{5}{7} \qquad\qquad
	8\dfrac{2}{7} \qquad\qquad
	4\dfrac{1}{10}  \]
}
\newpage
\eks[2]{
	Skriv om brøken $ \frac{17}{3} $ til et blandet tall.
	
	\sv 
	Vi legger merke til at telleren er 17 og nevneren er 3. Det største heltallet vi kan gange med 3 uten at produktet blir større enn 17, er 5. Altså kan vi skrive
	\alg{
		\frac{17}{\colb{3}}&=\frac{5\cdot\colb{3}+2}{\colb{3}}\br
		&= \frac{5\cdot\colb{\cancel{3}}}{\colb{\cancel{3}}}+\frac{2}{\colb{3}} \\
		&=5+\frac{2}{3}
	}
}
\eks[3]{
	Skriv om det blandede tallet $ 3+\frac{4}{5} $ til en brøk.
	
	\sv
	
	Vi har at $ 3=\frac{3}{1} $, altså er
	\[ 3+\frac{4}{5}=\frac{3}{1}+\frac{4}{5} \]
	Videre har vi at\footnote{Se \rref{braduliknemn}.}
	\alg{
		\frac{3}{1}+\frac{4}{5}&= \frac{3\cdot5}{1\cdot5}+\frac{4}{5} \br
		&= \frac{15}{5}+\frac{4}{5}\br
		&= \frac{19}{5}
	}
}

\end{document}


